\chapter{Mesurer Lightroom}

\section{Le protocole}
\subsection{Des paires d'exports comparables}
Chaque scène est exportée de Lightroom en TIFF 16 bits sRGB, une fois à l'état
neutre et une fois par valeur testée du curseur : $-100$, $-50$, $-25$, $+25$,
$+50$, $+100$. Les préréglages utilisés sont \emph{autonomes} : chacun fixe
tout l'état de développement (profil Camera Neutral, balance des blancs Prise
de vue, corrections optiques, bruit et netteté désactivés, tous les autres
curseurs à zéro) et ne change que le curseur étudié. Appliqué à n'importe
quelle photo, un préréglage donne donc un état connu, quels que soient ses
réglages précédents.

L'analyseur ne se fie pas au nom des fichiers. Il lit le XMP embarqué dans
chaque TIFF, y trouve la valeur du curseur, et associe l'export au neutre dont
\emph{tous les autres réglages} sont identiques. Une association ambiguë, un
état manquant ou dupliqué arrête l'analyse.

\subsection{Comparer un pixel à lui-même}
Dans une paire, le neutre et la variante ont la même géométrie : le pixel
$(x,y)$ de l'un correspond exactement au pixel $(x,y)$ de l'autre. On peut
donc mesurer, pour chaque pixel, l'effet du curseur :
\begin{equation}
\gamma = \ln\frac{C_{\text{variante}}}{C_{\text{neutre}}},\qquad
\lambda = \ln\frac{L_{\text{variante}}}{L_{\text{neutre}}},\qquad
\Delta h = h_{\text{variante}}-h_{\text{neutre}}\ (\text{ramené à }]-180^\circ,180^\circ]).
\end{equation}
Le logarithme rend les gains additifs et symétriques : doubler la chroma donne
$\gamma=\ln 2\approx0{,}69$, la diviser par deux donne $-0{,}69$.

\subsection{Échantillonnage et exclusions}
Un pixel sur 8 dans chaque direction est lu, puis 40\,000 à 60\,000 pixels par
état sont tirés au hasard avec une graine fixe (reproductible). Sont exclus :
\begin{itemize}
\item les pixels dont un canal est écrêté ($\le0{,}002$ ou $\ge0{,}995$ en
valeur encodée) dans l'un des deux exports : leur valeur ne dit plus rien de
l'effet du curseur ;
\item les pixels très sombres ($L\le0{,}08$), dominés par le bruit ;
\item pour les tables, les quasi-gris ($C\le0{,}01$), traités à part
(section~\ref{sec:quasigris}).
\end{itemize}

\resultat{Comparer chaque export à son propre neutre, et non à \IPR{},
mesure le comportement du curseur indépendamment du rendu de base.}

\section{Quelle variable explique la réponse ?}
On range les pixels en 20 classes d'une variable candidate, on retire de
chaque pixel la médiane de sa classe, et on mesure la variance restante de
$\gamma$. Plus elle est basse, mieux la variable explique la réponse.

\begin{center}
\begin{tabular}{lrrrr}
\toprule
Valeur & variance totale & $C/L$ Oklab & $C$ Oklab & $s$ ProPhoto\\
\midrule
$-100$ & 0,7465 & 0,5841 & \textbf{0,2031} & 0,5246\\
$-50$ & 0,0115 & 0,0078 & \textbf{0,0037} & 0,0054\\
$+50$ & 0,0184 & 0,0130 & 0,0183 & \textbf{0,0130}\\
$+100$ & 0,1355 & \textbf{0,0715} & 0,1344 & 0,0759\\
\bottomrule
\end{tabular}
\end{center}

Au positif, la saturation relative ($s$ ou $C/L$) explique bien plus que la
chroma absolue. Au négatif, la chroma absolue explique mieux, mais aucune
variable seule ne suffit.

On ajoute ensuite des variables, en croisant les classes. La part de variance
expliquée devient :

\begin{center}
\begin{tabular}{lrrrrr}
\toprule
Valeur & $s$ & $s$+$L$ & $s$+$L$+$h$ & $s$+scène & $s$+$L$+$h$+scène\\
\midrule
$-100$ & 73\,\% & 79\,\% & 85\,\% & 80\,\% & 88\,\%\\
$+25$ & 25\,\% & 52\,\% & 93\,\% & 56\,\% & 95\,\%\\
$+50$ & 29\,\% & 59\,\% & 94\,\% & 58\,\% & 96\,\%\\
$+100$ & 44\,\% & 77\,\% & 95\,\% & 67\,\% & 97\,\%\\
\bottomrule
\end{tabular}
\end{center}

Saturation, clarté et teinte expliquent 93 à 95\,\% de la variance au positif.
Connaître en plus l'image d'où vient le pixel n'apporte que 2\,\%.

\resultat{La vibrance de Lightroom est une fonction fixe du pixel, pas une
adaptation à l'image. Les grandes différences constatées d'une scène à
l'autre viennent de leur contenu : une scène riche en tons moyens peu saturés
réagit davantage.}

\section{La forme de la réponse}
\subsection{Le gain selon la saturation et la clarté}
Gain médian de chroma à $+100$, selon la saturation $s$ (lignes) et la clarté
$L$ (colonnes) du pixel neutre :

\begin{center}
\begin{tabular}{lrrrrrr}
\toprule
$s$ \textbackslash{} $L$ & 0,19 & 0,38 & 0,53 & 0,68 & 0,82 & 0,96\\
\midrule
0,00--0,10 & -- & 2,80 & 3,73 & 3,07 & 2,06 & 1,72\\
0,10--0,20 & 1,53 & 2,89 & 3,70 & 3,11 & 2,13 & 1,82\\
0,20--0,30 & 1,37 & 2,55 & 3,35 & 2,88 & 2,20 & 1,89\\
0,30--0,45 & 1,24 & 2,00 & 2,42 & 2,44 & 2,13 & 1,85\\
0,45--0,60 & 1,22 & 1,67 & 1,74 & 1,83 & 1,92 & 1,79\\
0,60--0,80 & 1,19 & 1,38 & 1,42 & 1,45 & 1,53 & --\\
0,80--1,00 & 1,13 & 1,22 & 1,31 & -- & -- & --\\
\bottomrule
\end{tabular}
\end{center}

Trois tendances se lisent :
\begin{itemize}
\item le gain décroît quand la saturation augmente : les couleurs déjà vives
sont protégées, c'est la définition même de la vibrance ;
\item il culmine sur les tons moyens ($L\approx0{,}5$) et retombe fortement
dans les ombres et les hautes lumières ;
\item la dépendance à $L$ est aussi forte que celle à $s$. L'ancien modèle, qui
l'ignorait, ne pouvait pas reproduire Lightroom.
\end{itemize}

\subsection{Une réponse progressive au curseur}
Au maximum du gain, la chroma est multipliée par 1,18 à $+25$, 1,42 à $+50$ et
2,49 à $+100$. En logarithme, $0{,}16$, $0{,}35$ et $0{,}91$ : la réponse
croît plus vite que le curseur, à peu près comme $V^{1{,}3}$. Le curseur est
doux au début et s'accélère vers $+100$.

\subsection{Le côté négatif}
À $-25$ et $-50$, le gain est presque uniforme (0,70 puis 0,45) jusqu'à
$s\approx0{,}5$ ; au-delà, les couleurs saturées sont protégées. À $-100$, les
couleurs faibles sont presque entièrement désaturées (gain 0,03 à 0,05), les
plus saturées gardent jusqu'à 0,35 de leur chroma.

\subsection{La clarté change aussi}
Lightroom ne garde pas la clarté constante. Pour les couleurs saturées, à
$-100$, $L$ est multipliée par 0,88 en médiane : désaturer assombrit. Ces
gains de clarté sont mesurés et reproduits comme ceux de la chroma. Ils
dépendent eux aussi de la couleur : certains bleus s'éclaircissent légèrement
en désaturant.

\subsection{Les quasi-gris}
\label{sec:quasigris}
Les pixels presque gris ($C<0{,}01$), souvent du bruit chromatique, reçoivent un
traitement particulier. À $+100$, Lightroom ne multiplie leur chroma que par
1,2 environ, quand leur classe de saturation prévoirait 1,55. À $-100$, au
contraire, il les désature davantage. La transition se fait entre $C=0{,}004$
et $C=0{,}012$.

\subsection{La teinte tourne}
Lightroom fait tourner les teintes en même temps qu'il les renforce. Rotation
médiane par famille de teinte, sur 13 scènes :

\begin{center}
\begin{tabular}{lrrr}
\toprule
Famille (teinte Oklab) & $+100$ & $+50$ & $-50$\\
\midrule
rouge (0°) & $+1{,}4^\circ$ & $+0{,}3^\circ$ & $-1{,}5^\circ$\\
orange (30°) & $+3{,}3^\circ$ & $+0{,}9^\circ$ & $-5{,}2^\circ$\\
jaune-orangé (60°) & $+4{,}6^\circ$ & $+2{,}0^\circ$ & $-3{,}6^\circ$\\
vert (150°) & $+0{,}1^\circ$ & $0{,}0^\circ$ & $0{,}0^\circ$\\
cyan (210°) & $-3{,}6^\circ$ & $-2{,}5^\circ$ & $+3{,}9^\circ$\\
bleu-cyan (240°) & $-7{,}7^\circ$ & $-2{,}5^\circ$ & $+5{,}4^\circ$\\
bleu (270°) & $-5{,}1^\circ$ & $-1{,}5^\circ$ & $+1{,}6^\circ$\\
\bottomrule
\end{tabular}
\end{center}

En renforçant, les bleus partent vers le cyan, les oranges vers le jaune ; en
désaturant, le mouvement s'inverse. Le chapitre 5 explique comment le modèle
le reproduit.

\resultat{La vibrance de Lightroom agit sur trois grandeurs à la fois : la
chroma, la clarté et la teinte. Toutes trois dépendent de la saturation, de la
clarté et de la teinte du pixel de départ.}
